\chapter*{Kurzfassung}
\label{chap:abstract}

Embedding-basierte Vektorsuchsysteme bilden das Rückgrat moderner Informationsrückgewinnung und agentenbasierter KI-Architekturen. Mit wachsenden Dokumentkorpora steigen Speicherbedarf und Betriebskosten hochdimensionaler Float32-Vektoren jedoch drastisch an. Die nachträgliche Quantisierung auf wenige Bits pro Dimension verspricht erhebliche Speichervorteile, führt jedoch unweigerlich zu Informationsverlusten. 

In dieser Arbeit wird der Einfluss von Quantisierung und Dimensionsreduktion auf die Retrieval-Qualität und den Ressourcenbedarf systematisch und kausal untersucht. Anhand eines kontrollierten Sechs-Phasen-Experiments auf Basis des Embedding-Modells \texttt{bge-large-en-v1.5} und dreier heterogener BEIR-Datensätze (SciFact, FiQA, TREC-COVID) wird die Erhaltung geometrischer Raumeigenschaften, Distanzkorrelationen, Nachbarschaften und Retrieval-Metriken (NDCG@10, Recall@100) über verschiedene Bittiefen (1 bis 8~Bit) sowie PCA-Zielstufen evaluiert. Verglichen werden naive Quantisierungsansätze mit dem rotationsbasierten Verfahren TurboQuant.

Die Ergebnisse zeigen, dass eine Kompression bis 4-Bit praktisch verlustfrei bleibt und über 99{,}5\,\% der Float32-Retrieval-Qualität bei einer Speicherreduktion um das 7{,}8-Fache wahrt. Für kostenoptimierte Cloud-Systeme bietet 2-Bit bei voller Dimension (1024d) den besten Kompromiss mit 98{,}1\,\% erhaltener Qualität und 15-facher Speicherersparnis, was unter den gewählten Annahmen einer modellierten AWS-Infrastruktur 93{,}7\,\% der monatlichen Instanzkosten einspart. Selbst 1-Bit-Konfigurationen behalten eine hohe Trustworthiness und eignen sich für zweistufige Filterstrukturen sowie lokale On-Device-Szenarien mit 100\,000 bis 500\,000 Vektoren unterhalb von 50~MiB Arbeitsspeicher. Die geometrische Analyse belegt zudem, dass TurboQuant durch orthogonale Rotation das bei nativer Binarisierung auftretende Verdünnungsparadox unter Dimensionsreduktion erfolgreich verhindert.

\vspace{0.5cm}
\noindent\textbf{Schlagwörter:} Vektor-Embeddings, Quantisierung, TurboQuant, Semantische Suche, Information Retrieval, Matryoshka-Reduktion, Skalierbarkeit